\daytitle{12}{万有引力与天体运动}

\begin{learningbox}{本节导读}
万有引力定律是牛顿最伟大的成就之一——它统一了“天上的力”和“地上的力”。
本节学习万有引力定律、天体运动的基本规律、三个宇宙速度，并
建立用万有引力解决天体问题的标准方法。
\end{learningbox}

\textbf{本节读题路线：} 先判断研究对象是在地面附近、绕天体做圆周运动，还是变轨过程；圆轨道优先写“万有引力提供向心力”，地面附近先由 $mg=GMm/R^2$ 得到 $GM=gR^2$。若题目涉及发射、升轨或逃逸，就要从速度公式扩展到机械能变化。

\section*{一、学习目标}

\begin{enumerate}
  \item 能写出万有引力定律 $F = G\dfrac{Mm}{r^2}$，并说明 $G$ 的物理意义。
  \item 能由 $mg=GMm/R^2$ 推出并使用黄金代换 $GM = gR^2$。
  \item 能由“万有引力提供向心力”写出天体运动基本方程 $G\dfrac{Mm}{r^2} = m\dfrac{v^2}{r}$。
  \item 能区分三个宇宙速度的物理含义，并会计算第一宇宙速度。
  \item 能列出地球静止卫星的轨道条件，并据此求轨道半径。
  \item 能用图像和语言概括开普勒三定律的初步内容。
\end{enumerate}
\section*{二、初高中衔接}

初中阶段只定性接触过“万有引力”这个概念，知道苹果落地是因为地球吸引。

高中把万有引力定量化，并揭示了一个深刻结论：\textbf{月球绕地球转和苹果落地都由同一种相互作用支配}——万有引力。牛顿的贡献在于：把天上的规律和地上的规律统一起来了。

\textbf{科学发展视角：} 现代导航卫星、气象卫星、深空探测器和空间站轨道设计，都离不开万有引力和圆周运动。学习本节时不要只把公式当作计算工具，还要形成一个观念：航天问题首先是“引力提供向心力”的模型问题。

\section*{三、核心概念}

\begin{conceptbox}{万有引力定律}

\[
\boxed{F = G\frac{Mm}{r^2}}
\]

\begin{itemize}
  \item $G = 6.67 \times 10^{-11}\ \mathrm{N \cdot m^2 / kg^2}$：引力常量（卡文迪许扭秤实验测得）。
  \item $M$、$m$：两个天体的质量。
  \item $r$：两个天体中心的距离。
  \item 方向：沿两物体中心的连线，相互吸引。
\end{itemize}

\begin{center}
% TikZ 图：万有引力定律示例
% Day 12：万有引力与天体运动
\begin{tikzpicture}[scale=0.8,>=Latex]
  % 两个天体
  \filldraw[blue!40,draw=blue] (0,0) circle (0.8);
  \node at (0,0) {$M$};
  \filldraw[red!40,draw=red] (4,0) circle (0.5);
  \node at (4,0) {$m$};
  % 中心距
  \draw[dashed] (0,0) -- (0,1.15);
  \draw[dashed] (4,0) -- (4,1.15);
  \draw[<->] (0,1.0) -- (4,1.0) node[midway,above] {$r$};
  % 引力：分别画在连线两侧，避免标签重叠
  \draw[->,thick,red] (0.85,-0.22) -- (1.85,-0.22) node[midway,below] {$F$};
  \draw[->,thick,red] (3.45,-0.22) -- (2.45,-0.22) node[midway,below] {$F'$};
  % 公式
  \node at (2,-1.5) {$F = G\frac{Mm}{r^2}$};
\end{tikzpicture}


\figcaption{万有引力：$F = F' = G\dfrac{Mm}{r^2}$，方向沿连线}\label{fig:gravity-law}
\end{center}

\textbf{适用提醒：}
\begin{enumerate}
  \item 万有引力定律直接适用于质点间的相互作用。对于质量呈球对称分布的球体，研究球体外部的引力时，可等效为质量集中在球心的质点；一般形状的物体不能无条件这样处理。
  \item $G$ 非常小，所以日常物体彼此之间的万有引力通常可忽略；但地球质量巨大，地球对物体的引力就是日常重力的来源。
\end{enumerate}

\end{conceptbox}

\begin{conceptbox}{重力与万有引力的关系}

\begin{center}
% TikZ 图：地球表面重力与万有引力
% Day 12：万有引力与天体运动
\begin{tikzpicture}[scale=0.7,>=Latex]
  % 地球
  \shadedraw[ball color=blue!20,draw=blue] (0,0) circle (2);
  \node at (0,0) {地球};
  % 表面物体
  \filldraw[red] (45:2) circle (3pt);
  \node[above right] at (45:2) {$m$};
  % 万有引力（指向地心）
  \draw[->,thick,red] (45:2) -- ++(225:1.5) node[midway,above right] {$F_{\text{引}}$};
  % 重力（竖直向下——近似指向地心）
  \draw[->,thick,blue] (45:2) -- ++(225:1.0) node[below left] {$G = mg$};
  % 自转轴
  \draw[dashed] (0,-2.8) -- (0,2.8);
  \node[above] at (0,2.8) {$N$（北极）};
  \node[below] at (0,-2.8) {$S$（南极）};
\end{tikzpicture}


\figcaption{地面附近物体：万有引力 $\approx$ 重力（忽略自转）}\label{fig:gravity-surface}
\end{center}

\textbf{黄金代换：}
\[
\boxed{GM = gR^2}
\]
其中 $M$ 为地球质量，$R$ 为地球半径，$g$ 为地表重力加速度。

\begin{itemize}
  \item 在地球表面，$mg \approx G\dfrac{Mm}{R^2}$，所以 $g = \dfrac{GM}{R^2}$。
  \item 在高度 $h$ 处：$g' = \dfrac{GM}{(R+h)^2} = g\left(\dfrac{R}{R+h}\right)^2$。
  \item 在地球外部，离地心越远，重力加速度越小。
\end{itemize}

\end{conceptbox}

\begin{conceptbox}{考虑地球自转时：引力、视重与向心力}

\vspace{0.2em}
前面把 $mg$ 近似看作万有引力，是\textbf{忽略地球自转}的结果。若考虑地球自转，地面物体随地球绕自转轴做圆周运动，物体所受合力必须指向自转轴。除极地外，秤的示数（视重）会略小于万有引力。

这里采用理想球形地球模型：水平秤的法线沿当地半径方向，秤的示数对应支持力 $N$ 的大小。真实的万有引力 $F_g$ 始终指向地心；在任意纬度处，地面还可能对物体产生沿切线方向的静摩擦力。支持力和静摩擦力合称地面对物体的接触力。

\vspace{0.6em}
\begin{center}
% TikZ 图：考虑地球自转时的受力
% Day 12：万有引力与天体运动
\begin{tikzpicture}[scale=0.78,>=Latex]
  \tikzset{
    earth/.style={draw=blue!70,thick,fill=blue!5},
    grav/.style={->,red,thick},
    normal/.style={->,orange,thick},
    centrip/.style={->,blue,thick},
    note/.style={font=\small,align=center}
  }

  % 极地
  \begin{scope}[xshift=-4.75cm]
    \coordinate (O) at (0,0);
    \coordinate (P) at (90:1.18);
    \draw[earth] (O) circle (1.18);
    \draw[dashed,gray] (0,-1.45) -- (0,1.45);
    \node[above] at (0,1.62) {极地};
    \filldraw[black] (O) circle (1.3pt);
    \filldraw[red] (P) circle (2.2pt);
    \draw[grav] (P) -- ++(270:0.66);
    \draw[normal] (P) -- ++(90:0.54);
    \node[red,left] at (-0.18,0.86) {$F_{\text{引}}$};
    \node[orange,right] at (0.12,1.48) {$N$};
    \node[note] at (0,-1.70) {极地：$\rho=0$};
  \end{scope}

  % 赤道
  \begin{scope}
    \coordinate (O) at (0,0);
    \coordinate (P) at (0:1.18);
    \draw[earth] (O) circle (1.18);
    \draw[dashed,gray] (0,-1.45) -- (0,1.45);
    \draw[dashed,gray] (-1.42,0) -- (1.42,0);
    \node[above] at (0,1.62) {赤道};
    \filldraw[black] (O) circle (1.3pt);
    \filldraw[red] (P) circle (2.2pt);
    \draw[grav] (P) -- ++(180:0.72);
    \draw[normal] (P) -- ++(0:0.55);
    \draw[centrip] ($(P)+(0,-0.30)$) -- ++(180:0.62);
    \node[red] at (0.55,0.33) {$F_{\text{引}}$};
    \node[orange] at (1.88,0.16) {$N$};
    \node[blue] at (0.63,-0.67) {$a_c$};
    \node[note] at (0,-1.70) {赤道：$\rho=R$};
  \end{scope}

  % 任意纬度
  \begin{scope}[xshift=4.75cm]
    \coordinate (O) at (0,0);
    \coordinate (P) at (38:1.18);
    \coordinate (H) at ({1.18*cos(38)},0);
    \draw[earth] (O) circle (1.18);
    \draw[dashed,gray] (0,-1.45) -- (0,1.45);
    \draw[dashed,gray] (-1.42,0) -- (1.42,0);
    \draw[dashed,gray] (O) -- (P);
    \draw[dashed,gray] (P) -- (H);
    \node[above] at (0,1.62) {纬度 $\varphi$};
    \filldraw[black] (O) circle (1.3pt);
    \filldraw[red] (P) circle (2.2pt);
    \draw[grav] (P) -- ++(218:0.66);
    \draw[normal] (P) -- ++(38:0.56);
    \draw[centrip] ($(P)+(0,0.22)$) -- ++(180:0.66);
    \draw[<->,thin] (0,-1.38) -- ($(H)+(0,-1.38)$);
    \node[red] at (0.43,0.45) {$F_{\text{引}}$};
    \node[orange] at (1.62,1.18) {$N$};
    \node[blue] at (0.50,1.20) {$a_c$};
    \node[note] at (0,-1.73) {纬度 $\varphi$：$\rho=R\cos\varphi$};
  \end{scope}
\end{tikzpicture}


\figcaption{考虑地球自转：万有引力与地面接触力的合力提供向心力}\label{fig:earth-rotation-weight}
\end{center}

设地球角速度为 $\omega$，纬度为 $\varphi$，地球半径为 $R$。

\begin{itemize}
  \item \textbf{极地：} 到自转轴距离为 $0$，不需要因自转产生的向心力，
  \[
  N_{\text{极}}=G\frac{Mm}{R^2}.
  \]
  \item \textbf{赤道：} 到自转轴距离为 $R$，向心力最大，
  \[
  G\frac{Mm}{R^2}-N_{\text{赤}}=m\omega^2R,
  \qquad
  N_{\text{赤}}=m\left(\frac{GM}{R^2}-\omega^2R\right).
  \]
  \item \textbf{纬度 $\varphi$ 处：} 到自转轴距离为 $R\cos\varphi$，向心加速度大小为 $\omega^2R\cos\varphi$，方向指向自转轴。把加速度分解到当地半径方向和切线方向，可得
  \[
  \boxed{N_{\varphi}=m\left(\frac{GM}{R^2}-\omega^2R\cos^2\varphi\right)},
  \qquad
  \boxed{f_{\varphi}=m\omega^2R\sin\varphi\cos\varphi}.
  \]
  其中静摩擦力沿当地切线指向较高纬度一侧。真实地球并非理想球体，精确的铅垂方向还要考虑地球形状和重力场分布；高中阶段通常只要求掌握极地、赤道及上述球形模型。
\end{itemize}

\vspace*{0.35em}
\textbf{认知提醒：} 万有引力是真实力，始终指向地心；向心力不是新力，而是合力指向自转轴的效果。水平秤的读数对应支持力大小，极地和赤道处各力共线；任意纬度必须保留矢量方向，不能把“万有引力减支持力”直接写成标量向心力方程。高中多数题默认忽略地球自转，因此直接用 $mg\approx G\dfrac{Mm}{R^2}$。

\end{conceptbox}

\begin{conceptbox}{天体运动的基本方程}

卫星绕地球做匀速圆周运动，且地球质量远大于卫星质量时，可近似认为地球中心不动，万有引力提供卫星的向心力：
\[
\boxed{G\frac{Mm}{r^2} = m\frac{v^2}{r} = m\omega^2 r = m\frac{4\pi^2}{T^2} r}
\]

由此可导出卫星运行参数：
\[
v = \sqrt{\frac{GM}{r}},\qquad
\omega = \sqrt{\frac{GM}{r^3}},\qquad
T = 2\pi\sqrt{\frac{r^3}{GM}}
\]

\textbf{规律：}
\begin{itemize}
  \item 轨道半径越大，速度越小、角速度越小、周期越大。
  \item 在地球外部的稳定圆轨道中，近地圆轨道的环绕速度最大（第一宇宙速度）；地球静止卫星的周期约为 $24\ \mathrm{h}$。
\end{itemize}

\begin{center}
% TikZ 图：万有引力与卫星轨道
% Day 12：万有引力与天体运动
\begin{tikzpicture}[scale=0.7,>=Latex]
  % 地球
  \filldraw[blue!30,draw=blue] (0,0) circle (1);
  \node at (0,0) {地球};
  % 卫星轨道
  \draw[dashed,thick] (0,0) circle (3.5);
  % 卫星
  \filldraw[red] (30:3.5) circle (3pt);
  \draw[->,thick,blue] (30:3.5) -- ++(120:1.2) node[above] {$v$};
  \draw[->,thick,red] (30:3.5) -- ++(210:1.2) node[below] {$F$};
  % 半径
  \draw[thick] (0,0) -- (30:3.5);
  \node at (15:2.2) {$r$};
\end{tikzpicture}


\figcaption{卫星绕地球：万有引力 $F$ 提供向心力}\label{fig:satellite}
\end{center}

\end{conceptbox}

\section*{四、公式与适用条件}

\textbf{公式使用前检查：}
\begin{enumerate}
  \item 先分清 $r$ 是两质心距离或轨道半径，不是离地高度 $h$。
  \item 判断是否可把天体看作质点；若用球心等效，需确认天体质量呈球对称分布且研究点在球体外部。
  \item 圆轨道题先写 $r=R+h$，再用万有引力提供向心力；椭圆轨道不能直接套匀速圆周公式。
  \item 地表附近用 $GM=gR^2$ 时，要说明忽略地球自转和高度变化。
\end{enumerate}

\begin{center}
\begin{tabular}{lll}
\toprule
公式 & 物理意义 & 适用条件 \\
\midrule
$F = G\dfrac{Mm}{r^2}$ & 万有引力 & 质点；球对称天体外部 \\
$GM = gR^2$ & 黄金代换 & 地表附近，忽略地球自转等修正 \\
$G\dfrac{Mm}{r^2} = m\dfrac{v^2}{r}$ & 天体运动基本方程 & 匀速圆周运动 \\
$T = 2\pi\sqrt{\dfrac{r^3}{GM}}$ & 卫星周期 & 绕中心天体做圆周运动 \\
\bottomrule
\end{tabular}
\end{center}

\begin{learningbox}{天体题先判三类入口}
万有引力题的入口并不只有一个。读题时先判断它属于哪一类：
\begin{enumerate}
  \item \textbf{地表附近：} 物体在天体表面或低高度附近，常用 $mg\approx G\dfrac{Mm}{R^2}$ 和 $GM=gR^2$，但要记得这里默认忽略自转等修正。
  \item \textbf{稳定圆轨道：} 卫星绕中心天体做匀速圆周运动，且中心天体质量远大于卫星质量时，核心方程是 $G\dfrac{Mm}{r^2}=m\dfrac{v^2}{r}$，这里的 $r$ 是到中心天体质心的距离。
  \item \textbf{发射或变轨：} 速度公式只描述稳定圆轨道；升轨、逃逸、转移椭圆要再看机械能是否增加，不能只比较圆轨道速度大小。
\end{enumerate}
先判入口，再代公式，能避免把“高度 $h$”“轨道半径 $r$”“天体半径 $R$”混在一起。
\end{learningbox}

\section*{五、典型模型}

\subsection*{模型 1：近地卫星与第一宇宙速度}

近地卫星：轨道半径 $\approx R$（地球半径），圆轨道环绕速度最大：
\[
G\frac{Mm}{R^2} = m\frac{v^2}{R},\qquad v = \sqrt{\frac{GM}{R}} = \sqrt{gR}
\]
第一宇宙速度
\[
v_1 = \sqrt{gR} = \sqrt{9.8 \times 6.4 \times 10^6}
\approx 7.9\ \mathrm{km/s}.
\]
它既是近地圆轨道的环绕速度，也是理想条件下从地面水平发射、使物体不再落回地面的最小速度。这里忽略空气阻力、地球自转和地形，并把地面附近看作可供卫星运行的理想圆轨道。

\textbf{第二、第三宇宙速度：}

\begin{center}
% TikZ 图：三个宇宙速度的两个空间尺度
% Day 12：万有引力与天体运动
\begin{tikzpicture}[scale=0.75,>=Latex,line cap=round,line join=round,
  every node/.style={font=\small}]
  % 地球尺度：第一、第二宇宙速度
  \begin{scope}[xshift=-3.55cm]
    \coordinate (E) at (0,0);
    \filldraw[blue!22,draw=blue!70,thick] (E) circle (0.62);
    \node at (E) {地球};
    \draw[dashed,gray] (E) circle (1.38);
    \draw[->,blue,very thick] (15:1.38) arc (15:105:1.38);
    \node[blue,align=center] at (-0.35,2.05) {$v_1$：近地环绕};
    \draw[->,orange!90!black,very thick]
      (35:0.63) .. controls (1.35,1.35) and (2.25,1.45) .. (2.90,0.80);
    \node[orange!90!black,align=center] at (2.15,1.55) {$v_2$：逃离地球};
    \node[font=\normalsize] at (0,-2.05) {地球尺度};
  \end{scope}

  % 太阳系尺度：第三宇宙速度
  \begin{scope}[xshift=3.65cm]
    \coordinate (S) at (0,0);
    \coordinate (E) at (2.25,0);
    \filldraw[yellow!60!orange,draw=orange!80!black,thick] (S) circle (0.55);
    \node at (S) {太阳};
    \draw[dashed,gray] (S) circle (2.25);
    \filldraw[blue!55,draw=blue!80!black] (E) circle (0.13);
    \node[below=4pt] at (E) {地球};
    \draw[->,red,very thick]
      ($(E)+(0.05,0.10)$) .. controls (2.30,0.95) and (3.05,1.65) .. (3.70,2.20);
    \node[red,align=left] at (3.30,2.72) {$v_3$：沿公转方向发射\\最终逃离太阳系};
    \draw[->,gray!75,thick] (125:2.25) arc (125:72:2.25);
    \node[gray!75] at (-0.05,2.67) {$v_{\mathrm E}$：地球公转方向};
    \node[font=\normalsize] at (0,-2.78) {太阳系尺度};
  \end{scope}
\end{tikzpicture}


\figcaption{三个宇宙速度对应不同尺度：近地环绕、逃离地球和逃离太阳系}\label{fig:cosmic-speeds}
\end{center}

\textbf{第二宇宙速度的来源：} 取无穷远处引力势能为零，质量为 $m$ 的物体在距地心 $r$ 处的引力势能为
\[
E_p=-\frac{GMm}{r}.
\]
物体恰好逃离地球时，到达无穷远处的速度趋于零。由机械能守恒，
\[
\frac12mv_2^2-\frac{GMm}{R}=0,
\]
所以
\[
\boxed{v_2=\sqrt{\frac{2GM}{R}}=\sqrt2\,v_1\approx11.2\ \mathrm{km/s}}.
\]

\textbf{第三宇宙速度的来源：} 地球绕太阳的公转速度约为
$v_{\mathrm E}=29.8\ \mathrm{km/s}$，在地球轨道处脱离太阳引力所需的日心速度为
$\sqrt2\,v_{\mathrm E}\approx42.1\ \mathrm{km/s}$。若沿地球公转方向发射，探测器脱离地球后相对地球仍需保留约
\[
v_{\infty}=(\sqrt2-1)v_{\mathrm E}\approx12.3\ \mathrm{km/s}
\]
的速度。采用分段处理的理想模型，有
\[
\boxed{v_3=\sqrt{v_2^2+v_{\infty}^2}\approx16.7\ \mathrm{km/s}}.
\]
这里忽略空气阻力、地球自转以及其他天体摄动；若发射方向不利，所需速度会更大。

\subsection*{模型 2：地球静止卫星}

若卫星要相对地面某点保持静止，它必须是\textbf{地球静止卫星}。它是同步卫星中的特殊类型：在赤道平面内沿圆轨道与地球同向转动，周期等于一个恒星日（约 $23\ \mathrm{h}\ 56\ \mathrm{min}$，高中计算常取 $24\ \mathrm{h}$）。
由 $T = 2\pi\sqrt{\dfrac{r^3}{GM}}$ 可算出轨道半径约 $4.2 \times 10^4\ \mathrm{km}$（距地面约 $3.6 \times 10^4\ \mathrm{km}$）。

\begin{center}
% TikZ 图：同步卫星模型
% Day 12：万有引力与天体运动
\begin{tikzpicture}[scale=0.82,>=Latex,line cap=round,line join=round,
  every node/.style={font=\small}]
  \coordinate (O) at (0,0);

  % Earth and equator.
  \filldraw[blue!18,draw=blue!70!black,thick] (O) circle (1.05);
  \draw[blue!70!black,thick] (-1.05,0) arc[start angle=180,end angle=360,x radius=1.05,y radius=0.28];
  \draw[blue!70!black,dashed] (1.05,0) arc[start angle=0,end angle=180,x radius=1.05,y radius=0.28];
  \node at (0,0.18) {地球};
  \node[blue!70!black] at (-1.80,0.28) {赤道平面};

  % Geostationary orbit.
  \draw[dashed,thick,gray!70] (O) circle (3.25);
  \draw[->,thick,gray!80] (120:3.25) arc[start angle=120,end angle=82,radius=3.25];
  \node[gray!80] at (1.15,3.30) {同向转动};

  % Satellite and fixed ground point.
  \filldraw[red!80!black] (0:3.25) circle (3pt);
  \node[red!80!black,right=3pt] at (0:3.25) {同步卫星};
  \filldraw[black] (0:1.05) circle (1.5pt);
  \node[below left=4pt] at (0:1.05) {地面点};
  \draw[densely dashed] (0:1.05) -- (0:3.25);

  % Radius and height.
  \draw[->,thick] (O) -- (0:3.25);
  \node[above=4pt] at (1.75,0) {$r$};
  \draw[orange!85!black,thick] (1.05,-0.20) -- (1.05,-0.38);
  \draw[orange!85!black,thick] (3.25,-0.20) -- (3.25,-0.38);
  \draw[<->,orange!85!black,thick] (1.05,-0.32) -- (3.25,-0.32);
  \node[orange!85!black,below=4pt] at (2.15,-0.32) {$h$};

  % Period relation.
  \node[align=center] at (0,-3.85)
    {$T_{\text{卫星}}=T_{\text{地球自转}}\approx24\ \mathrm{h}$\\
     $r\approx4.2\times10^4\ \mathrm{km}$，$h\approx3.6\times10^4\ \mathrm{km}$};
\end{tikzpicture}


\figcaption{地球静止卫星：圆形赤道轨道、同向转动且周期等于一个恒星日}\label{fig:geostationary-satellite}
\end{center}

\textbf{地球静止卫星的四个条件：}
\begin{enumerate}
  \item 轨道是圆轨道；
  \item 轨道平面是赤道平面；
  \item 运动方向与地球自转方向相同；
  \item 周期等于地球自转的恒星周期。
\end{enumerate}

满足这四个条件时，卫星相对地面静止，始终位于赤道上某一经度的上空，常用于通信、广播和气象观测。不同地球静止卫星可以位于不同经度上方，但轨道半径和周期相同。

\subsection*{模型 3：卫星变轨}

\textbf{先看圆轨道的能量。} 由万有引力提供向心力，
\[
\frac{GMm}{r^2}=\frac{mv^2}{r}
\quad\Longrightarrow\quad
E_k=\frac12mv^2=\frac{GMm}{2r}.
\]
取无穷远处引力势能为零，则
\[
E_p=-\frac{GMm}{r},
\qquad
\boxed{E=E_k+E_p=-\frac{GMm}{2r}}.
\]
轨道半径越大，圆轨道速度越小，但机械能越大（负值的绝对值越小）。

\textbf{再看两次变轨。} 从低圆轨道升到高圆轨道，发动机必须对卫星做正功。典型过程是：在低轨道的近地点沿速度方向第一次加速，进入转移椭圆；到达远地点后再沿速度方向加速，使轨道圆化到高轨道。

\begin{center}
\begin{tikzpicture}[scale=0.92,>=Stealth]
  \coordinate (O) at (0,0);
  \fill[blue!15] (O) circle (0.34);
  \node[font=\scriptsize] at (0,-0.58) {地球};
  \draw[thick] (O) circle (1.15);
  \draw[thick] (O) circle (2.15);
  \draw[orange!85!black,very thick] (-0.50,0) ellipse (1.65 and 1.57);
  \fill[red] (1.15,0) circle (0.055);
  \fill[red] (-2.15,0) circle (0.055);
  \draw[->,red,very thick] (1.15,0) -- (1.15,0.62);
  \node[font=\scriptsize,red,right] at (1.20,0.52) {近地点加速 $\Delta v_1$};
  \draw[->,red,very thick] (-2.15,0) -- (-2.15,-0.62);
  \node[font=\scriptsize,red,left] at (-2.20,-0.55) {远地点加速 $\Delta v_2$，圆化};
  \draw[->,blue,thick] (0.80,0.82) arc (46:76:1.15);
  \draw[->,blue,thick] (-1.01,1.90) arc (118:150:2.15);
  \node[font=\small] at (1.70,-1.00) {低圆轨道};
  \node[font=\small] at (-2.45,1.58) {高圆轨道};
  \node[font=\small,orange!85!black] at (0.55,1.62) {转移椭圆};
  \node[font=\small,align=center] at (0,-2.55) {升到更高圆轨道：近地点加速进入椭圆，远地点再加速圆化};
\end{tikzpicture}


\figcaption{卫星升轨：近地点加速进入转移椭圆，远地点再加速圆化}\label{fig:satellite-transfer-orbit}
\end{center}

\textbf{认知提醒：} “高轨道圆周速度更小”和“升轨需要加速”并不矛盾。前者比较的是两个稳定圆轨道上的速度；后者讨论的是变轨瞬间的速度增量和机械能变化。比较速度时必须说明是在圆轨道上，还是在转移椭圆的近地点、远地点。

\subsection*{模型 4：双星系统}

两颗星绕共同质心做圆周运动，距离 $L = r_1 + r_2$：

\begin{center}
% TikZ 图：双星系统
% Day 12：万有引力与天体运动
\begin{tikzpicture}[scale=0.98,>=Latex]
  \coordinate (O) at (0,0);
  \coordinate (A) at (-1.55,0);
  \coordinate (B) at (2.10,0);

  \draw[dashed,gray] (O) circle (1.55);
  \draw[dashed,gray] (O) circle (2.10);
  \filldraw[black] (O) circle (2pt);
  \node[below=5pt] at (O) {$O$};

  \draw[<->,thin] ($(A)+(0,1.55)$) -- ($(B)+(0,1.55)$)
    node[midway,above] {$L=r_1+r_2$};
  \draw[<->,thin] ($(A)+(0,-1.20)$) -- ($(O)+(0,-1.20)$)
    node[midway,below] {$r_1$};
  \draw[<->,thin] ($(O)+(0,-1.20)$) -- ($(B)+(0,-1.20)$)
    node[midway,below] {$r_2$};

  \filldraw[blue!35,draw=blue,thick] (A) circle (0.34);
  \node[blue,above left=3pt] at (A) {$M$};
  \filldraw[red!35,draw=red,thick] (B) circle (0.25);
  \node[red,above right=4pt] at (B) {$m$};

  \draw[->,blue,very thick] (145:1.55) arc (145:215:1.55);
  \node[blue,left] at (205:1.78) {$v_1$};
  \draw[->,red,very thick] (-35:2.10) arc (-35:35:2.10);
  \node[red,right] at (28:2.34) {$v_2$};

  \draw[->,purple,thick] ($(A)+(0,-0.28)$) -- ++(0:0.86);
  \draw[->,purple,thick] ($(B)+(0,-0.28)$) -- ++(180:0.86);
  \node[purple] at (-0.92,-0.70) {$F$};
  \node[purple] at (1.46,-0.70) {$F$};

  \node[font=\small,align=center] at (0,-1.88)
    {两星绕共同质心 $O$ 转动：角速度相同、周期相同，速度方向在连线两端相反。};
\end{tikzpicture}


\figcaption{双星系统：两星绕共同质心 $O$ 做同周期圆周运动}\label{fig:binary-star}
\end{center}

设两星质量分别为 $M$、$m$，到共同质心的距离分别为 $r_1$、$r_2$。质心位置满足
\[
Mr_1=mr_2,\qquad r_1+r_2=L.
\]
所以
\[
r_1=\frac{m}{M+m}L,\qquad r_2=\frac{M}{M+m}L.
\]

两星之间的万有引力大小相同，分别提供各自圆周运动的向心力：
\[
G\frac{Mm}{L^2} = M\omega^2 r_1 = m\omega^2 r_2
\]
代入 $r_1=\dfrac{m}{M+m}L$，得
\[
G\frac{Mm}{L^2}=M\omega^2\frac{m}{M+m}L
\]
化简：
\[
\boxed{\omega^2=\frac{G(M+m)}{L^3}},\qquad
\boxed{T=2\pi\sqrt{\frac{L^3}{G(M+m)}}}.
\]

\textbf{结论：} 两星角速度相同、周期相同；轨道半径与质量成反比：
\[
\frac{r_1}{r_2}=\frac{m}{M}.
\]
质量大的星离质心近，质量小的星离质心远。

\section*{六、例题精讲}

\begin{examplebox}{天体质量的计算}

\textbf{题意分析：} 已知地球半径 $R = 6.4 \times 10^6\ \mathrm{m}$，地表重力加速度 $g = 9.8\ \mathrm{m/s^2}$，引力常量 $G = 6.67 \times 10^{-11}\ \mathrm{N\cdot m^2/kg^2}$。求地球质量 $M$。

\textbf{物理模型：} 黄金代换。

\textbf{解题过程：}

由 $GM = gR^2$：
\[
M = \frac{gR^2}{G} = \frac{9.8 \times (6.4 \times 10^6)^2}{6.67 \times 10^{-11}}
\]
\[
= \frac{9.8 \times 4.1 \times 10^{13}}{6.67 \times 10^{-11}} \approx 6.0 \times 10^{24}\ \mathrm{kg}
\]

\textbf{方法总结：} 天体表面题先写 $mg\approx G\frac{Mm}{R^2}$。若已知 $g$ 和 $R$，并忽略自转等修正，就可用 $GM=gR^2$ 或进一步估算天体质量。

\end{examplebox}

\begin{examplebox}{卫星运行参数}

\textbf{题意分析：} 一颗卫星在距地面 $h = 400\ \mathrm{km}$ 的轨道上运行，$R = 6.4 \times 10^6\ \mathrm{m}$，$M = 6.0 \times 10^{24}\ \mathrm{kg}$，$G = 6.67 \times 10^{-11}\ \mathrm{N\cdot m^2/kg^2}$。求卫星的线速度、角速度和周期。

\textbf{物理模型：} 万有引力提供向心力。

\textbf{解题过程：}

轨道半径 $r = R + h = 6.4 \times 10^6 + 4.0 \times 10^5 = 6.8 \times 10^6\ \mathrm{m}$。

线速度：
\[
\begin{aligned}
v &= \sqrt{\frac{GM}{r}}
 = \sqrt{\frac{6.67 \times 10^{-11} \times 6.0 \times 10^{24}}{6.8 \times 10^6}} \\
&\approx \sqrt{5.88 \times 10^7}
 \approx 7.67 \times 10^3\ \mathrm{m/s}
 = 7.67\ \mathrm{km/s}
\end{aligned}
\]

角速度：
\[
\omega = \frac{v}{r} = \frac{7.67 \times 10^3}{6.8 \times 10^6} \approx 1.13 \times 10^{-3}\ \mathrm{rad/s}
\]

周期：
\[
T = \frac{2\pi}{\omega} = \frac{2 \times 3.14}{1.13 \times 10^{-3}} \approx 5.56 \times 10^3\ \mathrm{s} \approx 1.54\ \mathrm{h}
\]

\textbf{方法总结：} 卫星问题先把离地高度换成轨道半径 $r=R+h$，再写 $F_{\text{引}}=F_{\text{向}}$ 列方程；半径、速度和周期单位要先统一。

\end{examplebox}

\begin{examplebox}{地球静止卫星轨道半径}

\textbf{题意分析：} 地球静止卫星的周期取 $T=24\ \mathrm{h}$，地球质量 $M=6.0\times10^{24}\ \mathrm{kg}$，引力常量 $G=6.67\times10^{-11}\ \mathrm{N\cdot m^2/kg^2}$。求卫星的轨道半径，并估算离地高度。地球半径取 $R=6.4\times10^6\ \mathrm{m}$。

\textbf{物理模型：} 地球静止卫星沿赤道上空的圆轨道运行，万有引力提供向心力，周期由地球自转的恒星周期决定。

\begin{center}
% TikZ 图：同步卫星模型
% Day 12：万有引力与天体运动
\begin{tikzpicture}[scale=0.82,>=Latex,line cap=round,line join=round,
  every node/.style={font=\small}]
  \coordinate (O) at (0,0);

  % Earth and equator.
  \filldraw[blue!18,draw=blue!70!black,thick] (O) circle (1.05);
  \draw[blue!70!black,thick] (-1.05,0) arc[start angle=180,end angle=360,x radius=1.05,y radius=0.28];
  \draw[blue!70!black,dashed] (1.05,0) arc[start angle=0,end angle=180,x radius=1.05,y radius=0.28];
  \node at (0,0.18) {地球};
  \node[blue!70!black] at (-1.80,0.28) {赤道平面};

  % Geostationary orbit.
  \draw[dashed,thick,gray!70] (O) circle (3.25);
  \draw[->,thick,gray!80] (120:3.25) arc[start angle=120,end angle=82,radius=3.25];
  \node[gray!80] at (1.15,3.30) {同向转动};

  % Satellite and fixed ground point.
  \filldraw[red!80!black] (0:3.25) circle (3pt);
  \node[red!80!black,right=3pt] at (0:3.25) {同步卫星};
  \filldraw[black] (0:1.05) circle (1.5pt);
  \node[below left=4pt] at (0:1.05) {地面点};
  \draw[densely dashed] (0:1.05) -- (0:3.25);

  % Radius and height.
  \draw[->,thick] (O) -- (0:3.25);
  \node[above=4pt] at (1.75,0) {$r$};
  \draw[orange!85!black,thick] (1.05,-0.20) -- (1.05,-0.38);
  \draw[orange!85!black,thick] (3.25,-0.20) -- (3.25,-0.38);
  \draw[<->,orange!85!black,thick] (1.05,-0.32) -- (3.25,-0.32);
  \node[orange!85!black,below=4pt] at (2.15,-0.32) {$h$};

  % Period relation.
  \node[align=center] at (0,-3.85)
    {$T_{\text{卫星}}=T_{\text{地球自转}}\approx24\ \mathrm{h}$\\
     $r\approx4.2\times10^4\ \mathrm{km}$，$h\approx3.6\times10^4\ \mathrm{km}$};
\end{tikzpicture}

\end{center}

\textbf{解题过程：}

轨道周期满足
\[
T=2\pi\sqrt{\frac{r^3}{GM}}.
\]
整理得
\[
r=\sqrt[3]{GM\left(\frac{T}{2\pi}\right)^2}.
\]
把 $T=24\times3600=8.64\times10^4\ \mathrm{s}$ 代入：
\[
r\approx\sqrt[3]{6.67\times10^{-11}\times6.0\times10^{24}
\left(\frac{8.64\times10^4}{2\pi}\right)^2}
\approx4.2\times10^7\ \mathrm{m}.
\]
离地高度为
\[
h=r-R\approx4.2\times10^7-6.4\times10^6
\approx3.6\times10^7\ \mathrm{m}.
\]

\textbf{方法总结：} 地球静止卫星题先核对圆形赤道轨道、同向和同周期四项条件，再由周期求轨道半径 $r$，最后减去地球半径得到离地高度 $h$。不要把题中“高度”和“轨道半径”混用。
\end{examplebox}

\begin{examplebox}{双星系统的轨道半径与周期}

\textbf{题意分析：} 两颗恒星的质量分别为 $M=2.0\times10^{30}\ \mathrm{kg}$、
$m=1.0\times10^{30}\ \mathrm{kg}$，两星中心距离
$L=3.0\times10^{11}\ \mathrm{m}$，近似绕共同质心做圆周运动。求两星轨道半径之比，并估算公转周期。取
$G=6.67\times10^{-11}\ \mathrm{N\cdot m^2/kg^2}$。

\textbf{物理模型：} 双星绕共同质心做同角速度圆周运动。

\textbf{解题过程：}

由质心关系
\[
Mr_1=mr_2
\quad\Longrightarrow\quad
\frac{r_1}{r_2}=\frac{m}{M}=\frac12.
\]
又因为 $r_1+r_2=L$，所以
\[
r_1=1.0\times10^{11}\ \mathrm{m},
\qquad
r_2=2.0\times10^{11}\ \mathrm{m}.
\]
双星周期满足
\[
T=2\pi\sqrt{\frac{L^3}{G(M+m)}}.
\]
代入数据：
\[
T=2\pi\sqrt{\frac{(3.0\times10^{11})^3}
{6.67\times10^{-11}\times3.0\times10^{30}}}
\approx7.3\times10^7\ \mathrm{s}
\approx2.3\ \text{年}.
\]

\textbf{方法总结：} 双星题先用质心关系分配两个轨道半径，再用两星中心距离 $L$ 写万有引力；公式中的 $L$ 不能误写成某一颗星自己的轨道半径。

\end{examplebox}

\section*{七、变式训练}

\begin{enumerate}
  \item 已知月球半径 $R_{\text{月}} = 1.74 \times 10^6\ \mathrm{m}$，月球表面重力加速度 $g_{\text{月}} = 1.62\ \mathrm{m/s^2}$。取 $G=6.67\times10^{-11}\ \mathrm{N\cdot m^2/kg^2}$，求月球的质量。

  \item 一颗卫星在半径为 $r_1$ 的低圆轨道上运行，在某点沿速度方向短时加速后进入转移椭圆，并在远地点再次沿速度方向加速进入半径为 $r_2$ 的高圆轨道。比较第一次加速前后、第二次加速前后的速度和机械能变化，并说明为什么不能用“高轨道圆周速度更小”否定两次加速。

  \item 两颗卫星 $A$、$B$ 绕地球做匀速圆周运动，轨道半径之比 $r_A : r_B = 1 : 4$。求线速度之比和周期之比。
  \begin{center}
  % 变式训练：两颗卫星的圆轨道半径之比为 1:4
\begin{tikzpicture}[scale=0.72,>=Latex,line cap=round,line join=round]
  \coordinate (O) at (0,0);
  \coordinate (A) at (0:0.65);
  \coordinate (B) at (0:2.60);

  \filldraw[blue!15,draw=blue!65,thick] (O) circle (0.16);
  \node[below left=1pt,font=\scriptsize] at (O) {$O$};
  \draw[thick] (O) circle (0.65);
  \draw[thick] (O) circle (2.60);

  \filldraw[black] (A) circle (2pt);
  \node[below right=2pt] at (A) {$A$};
  \filldraw[black] (B) circle (2pt);
  \node[below right=2pt] at (B) {$B$};

  \draw[->,blue,thick] (A) -- ++(90:0.78);
  \node[blue,left] at ($(A)+(0,0.62)$) {$v_A$};
  \draw[->,red,thick] (B) -- ++(90:0.42);
  \node[red,right] at ($(B)+(0,0.36)$) {$v_B$};

  \node[font=\small] at (0,3.05) {$r_A:r_B=1:4$（按比例绘制）};
\end{tikzpicture}

  \end{center}

  \item 某双星系统中两星质量之比为 $M:m=3:1$，两星中心距离为 $L$。求它们到共同质心的距离之比，并写出该双星系统的周期表达式。
\end{enumerate}

\begin{learningbox}{变式训练参考结果}
\begin{enumerate}[itemsep=0.45em,topsep=0.3em,parsep=0pt]
  \item 由 $g_{\text{月}}=GM_{\text{月}}/R_{\text{月}}^2$，得
  \[
  M_{\text{月}}=\frac{g_{\text{月}}R_{\text{月}}^2}{G}
  \approx7.35\times10^{22}\ \mathrm{kg}.
  \]
  \item 两次短时加速都使卫星在加速点的速度增大、机械能增大：第一次由低圆轨道进入转移椭圆，第二次由转移椭圆进入高圆轨道。“高圆轨道速度较小”比较的是不同位置上的两条圆轨道；两次变轨比较的是同一位置加速前后的速度，二者并不矛盾。
  \item 由 $v\propto r^{-1/2}$、$T\propto r^{3/2}$，得
  $v_A:v_B=2:1$，$T_A:T_B=1:8$。
  \item 到共同质心的距离满足 $Mr_M=mr_m$，故
  $r_M:r_m=1:3$，且 $r_M=L/4$、$r_m=3L/4$。周期为
  \[
  T=2\pi\sqrt{\frac{L^3}{G(M+m)}}.
  \]
\end{enumerate}
\end{learningbox}

\section*{八、竞赛入门补充}

\begin{competitionbox}
\textbf{从开普勒三定律到牛顿万有引力}

在牛顿之前，开普勒已经从第谷的大量天文观测中总结出行星运动的经验规律：

\begin{enumerate}
  \item \textbf{椭圆轨道定律：} 行星绕太阳运动的轨道是椭圆，太阳在椭圆的一个焦点上。
  \item \textbf{面积定律：} 行星与太阳的连线在相等时间内扫过的面积相等（近地点速度快，远地点速度慢）。
  \item \textbf{周期定律：} 对绕同一中心天体运动的行星，$\dfrac{T^2}{a^3} = \text{常量}$，其中 $a$ 为椭圆轨道的半长轴。
\end{enumerate}

\begin{center}
% TikZ 图：开普勒椭圆轨道定律与面积定律
% Day 12：万有引力与天体运动
\begin{tikzpicture}[scale=0.76,>=Latex,line cap=round,line join=round,
  every node/.style={font=\small}]
  % 第一定律：太阳位于焦点
  \begin{scope}[xshift=-3.55cm]
    \coordinate (O) at (0,0);
    \coordinate (S) at (-1.56,0);
    \draw[thick] (O) ellipse (2.0 and 1.25);
    \filldraw[yellow!65!orange,draw=orange!80!black] (S) circle (0.18);
    \node[below=4pt] at (S) {太阳（焦点）};
    \filldraw[blue!55,draw=blue!80!black] (1.35,0.92) circle (0.10);
    \node[above right=2pt] at (1.35,0.92) {行星};
    \filldraw[black] (O) circle (1.2pt);
    \draw[<->,thin] (-2.0,-1.55) -- (2.0,-1.55);
    \node[below=2pt] at (0,-1.55) {$2a$};
    \node[font=\normalsize] at (0,1.80) {椭圆轨道定律};
  \end{scope}

  % 第二定律：相等时间扫过相等面积
  \begin{scope}[xshift=3.45cm]
    \coordinate (O) at (0,0);
    \coordinate (S) at (-1.56,0);
    \draw[thick] (O) ellipse (2.0 and 1.25);
    \filldraw[yellow!65!orange,draw=orange!80!black] (S) circle (0.18);
    \node[below=3pt] at (S) {太阳};

    % 近、远日点附近的两段扫掠区域（示意）
    \fill[red!22]
      (S) -- (-1.86,0.46)
      arc[start angle=158,end angle=202,x radius=2.0,y radius=1.25]
      -- cycle;
    \draw[red,thick] (S) -- (-1.86,0.46);
    \draw[red,thick] (S) -- (-1.86,-0.46);

    \fill[blue!20]
      (S) -- (1.997,-0.065)
      arc[start angle=-3,end angle=3,x radius=2.0,y radius=1.25]
      -- cycle;
    \draw[blue,thick] (S) -- (1.997,-0.065);
    \draw[blue,thick] (S) -- (1.997,0.065);

    \node[red] at (-1.40,0.76) {$\Delta A_1$};
    \node[blue] at (0.75,0.28) {$\Delta A_2$};
    \node[align=center] at (0,-1.75)
      {$\Delta t_1=\Delta t_2$ 时，$\Delta A_1=\Delta A_2$};
    \node[font=\normalsize] at (0,1.80) {面积定律};
  \end{scope}
\end{tikzpicture}


\figcaption{开普勒前两定律：太阳位于椭圆焦点，相等时间内扫过的面积相等}\label{fig:kepler-laws}
\end{center}

\textbf{牛顿怎样从圆周运动走向万有引力定律}

牛顿的关键思路不是先“猜一个公式”，而是把地面落体、月球绕地和开普勒定律放到同一个模型中。

\begin{center}
% TikZ 图：牛顿的月地检验
% Day 12：万有引力与天体运动
\begin{tikzpicture}[scale=0.78,>=Latex]
  \coordinate (E) at (0,0);
  \coordinate (Apple) at (0.86,0);
  \coordinate (Moon) at (4.95,0);

  \draw[dashed,gray] (E) circle (4.95);
  \filldraw[blue!20,draw=blue!70,thick] (E) circle (0.75);
  \node at (E) {地球};

  \filldraw[red] (Apple) circle (2.5pt);
  \node[above=5pt] at (Apple) {苹果};
  \draw[->,red,thick] (Apple) -- (0.28,0);
  \node[red,below=4pt] at (0.57,0) {$g$};

  \filldraw[gray!45,draw=gray,thick] (Moon) circle (0.23);
  \node[below right=4pt] at (Moon) {月球};
  \draw[->,blue,thick] (Moon) -- ++(90:0.82);
  \node[blue,above] at ($(Moon)+(0,0.82)$) {$v$};
  \draw[->,red,thick] (Moon) -- (3.82,0);
  \node[red,below=4pt] at (4.36,0) {$a_{\text{月}}$};

  \draw[thin,gray] (E) -- (Moon);
  \draw[<->,thin] (0,-1.05) -- (4.95,-1.05)
    node[midway,below] {$r\approx 60R$};

  \node[align=left,font=\small] at (7.45,0.45)
    {$a_{\text{月}}=\dfrac{4\pi^2r}{T^2}$\\[5pt]
     $a_{\text{月}}\approx \dfrac{g}{60^2}$\\[5pt]
     $\Rightarrow F\propto \dfrac{1}{r^2}$};
\end{tikzpicture}


\figcaption{牛顿的月地检验：月球所需向心加速度约为地表重力加速度的 $1/60^2$}\label{fig:newton-moon-test}
\end{center}

\textbf{第一步：月球绕地球需要向心加速度。}

把月球近似看成绕地球做圆周运动：
\[
a_{\text{月}}=\frac{4\pi^2r}{T^2}.
\]
代入月地距离约 $r\approx 60R$、月球周期约 $T\approx 27.3\ \text{天}$，可得
\[
a_{\text{月}}\approx 2.7\times 10^{-3}\ \mathrm{m/s^2}.
\]

\textbf{第二步：与地表重力加速度比较。}

地表重力加速度 $g\approx 9.8\ \mathrm{m/s^2}$，而
\[
\frac{g}{60^2}\approx \frac{9.8}{3600}\approx 2.7\times 10^{-3}\ \mathrm{m/s^2}.
\]
这说明：如果地球对物体的吸引随距离按平方反比减弱，那么地表的“苹果下落”和月球绕地的向心加速度可以由同一种力解释。

\textbf{第三步：用开普勒第三定律得到平方反比。}

对近似圆轨道，
\[
a_n=\frac{4\pi^2r}{T^2}.
\]
开普勒第三定律给出 $T^2\propto r^3$，所以
\[
a_n\propto \frac{r}{r^3}=\frac{1}{r^2}.
\]
由 $F=ma$，行星受到的中心力也满足
\[
F\propto \frac{m}{r^2}.
\]
再结合“相互作用力成对出现”和不同中心天体的比较，可得到中心天体质量也应成正比进入公式，于是形成
\[
\boxed{F=G\frac{Mm}{r^2}}.
\]
在牛顿理论中，对绕质量为 $M$ 的同一中心天体运动的椭圆轨道，可进一步证明
\[
\boxed{\frac{T^2}{a^3}=\frac{4\pi^2}{GM}}.
\]
这一步把开普勒总结出的经验常量，与中心天体质量和引力常量联系起来。

\textbf{例题：} 地球公转周期 $T_{\text{地}} = 1\ \text{年}$，轨道半径 $R_{\text{地}} = 1\ \text{AU}$。某小行星轨道半长轴 $a = 2.5\ \text{AU}$，求其公转周期。

\textbf{解：} 由 $\dfrac{T^2}{a^3} = \text{常量}$：
\[
\frac{T_{\text{行}}^2}{T_{\text{地}}^2} = \frac{a^3}{R_{\text{地}}^3},\qquad
T_{\text{行}} = T_{\text{地}} \cdot \sqrt{\frac{a^3}{R_{\text{地}}^3}} = 1 \times \sqrt{2.5^3} \approx 3.95\ \text{年}
\]

\textbf{竞赛提示：} 开普勒第三定律对绕同一中心天体的所有卫星/行星都成立，$T^2 \propto r^3$（椭圆轨道中为 $a^3$）。牛顿的高明之处在于：把这个经验规律转化成了“中心力按 $1/r^2$ 衰减”的动力学定律。
\end{competitionbox}

\section*{九、易错点}

\begin{mistakebox}
\small
\begin{enumerate}[nosep]
  \item \textbf{混淆 $r$ 的含义：} 先画地心、地表和卫星位置。地表问题的 $R$ 是地球半径，卫星问题的 $r$ 是地心到卫星的轨道半径。

  \item \textbf{误以为高轨道卫星速度更大：} 先用 $v=\sqrt{GM/r}$ 比较圆轨道速度；轨道半径越大速度越小，但总机械能更高。

  \item \textbf{把地球静止卫星理解成任意周期为 24 h 的卫星：} 判断时按四项检查：圆轨道、赤道平面、与地球自转同向、周期等于一个恒星日。仅有周期相同还不能相对地面静止。

  \item \textbf{忘记黄金代换的推导：} 需要消去 $GM$ 时，先从地表近似 $mg \approx G\frac{Mm}{R^2}$ 推出 $GM=gR^2$。

  \item \textbf{单位错误：} 代入前先把半径、高度和速度统一到国际单位。天体问题中 $R$ 通常是 $10^6\ \mathrm{m}$ 量级，$M$ 是 $10^{24}\ \mathrm{kg}$ 量级，计算时建议全程用科学计数法。

  \item \textbf{混淆三个宇宙速度：} 第一宇宙速度（$7.9\ \mathrm{km/s}$）是环绕速度，第二宇宙速度（$11.2\ \mathrm{km/s}$）是脱离地球引力束缚，第三宇宙速度（$16.7\ \mathrm{km/s}$）是脱离太阳系引力束缚。这些数值通常指从地面附近发射，并忽略空气阻力和地球自转影响时的理想值。
\end{enumerate}
\end{mistakebox}

\section*{十、今日小结}
\enlargethispage{4\baselineskip}

{\small
\begin{enumerate}[nosep]
  \item 万有引力 $F = G\dfrac{Mm}{r^2}$；地表附近忽略自转时可用黄金代换 $GM=gR^2$。
  \item 卫星方程：$G\dfrac{Mm}{r^2} = m\dfrac{v^2}{r}$ $\to$ $v = \sqrt{\dfrac{GM}{r}}$，$T = 2\pi\sqrt{\dfrac{r^3}{GM}}$。
  \item 轨道半径越大，$v$、$\omega$ 越小，$T$ 越大。
  \item 地球静止卫星沿圆形赤道轨道与地球同向转动，周期等于一个恒星日，轨道半径约 $4.2\times10^4\ \mathrm{km}$。
  \item 三个宇宙速度：$v_1\approx 7.9\ \mathrm{km/s}$，$v_2\approx 11.2\ \mathrm{km/s}$，$v_3\approx 16.7\ \mathrm{km/s}$。
  \item 卫星升轨需要增加机械能；高轨道圆周速度较小，但变轨瞬间常需要沿运动方向加速。
\end{enumerate}

\textbf{一分钟复盘：} 画一颗地球和一条圆轨道，标出 $r=R+h$，写出万有引力提供向心力的方程；再口述 $r$ 增大时 $v$、$T$、圆轨道机械能怎样变化，并说清地球静止卫星和三个宇宙速度各自对应的条件。
}
